% --- Archon physics formalization source begin ---
% archon:physics
% archon:covers PhyXMiniProblems/problem_phyx_mini_0684.lean
% archon:source-report reports/phyx_mini/problem_phyx_mini_0684.source.json

\chapter{Physics problem phyx\_mini\_0684}
\label{ch:PhyXMiniProblems_problem_phyx_mini_0684}

\paragraph{Problem source.}
\#\# Physical scenario

A landscape architect is planning an artificial waterfall in a city park. Water flowing at \textbackslash{}( 1.70 \textbackslash{}text\{ m/s\} \textbackslash{}) will leave the end of a horizontal channel at the top of a vertical wall \textbackslash{}( h = 2.35 \textbackslash{}text\{ m\} \textbackslash{}) high, and from there it will fall into a pool as shown in figure.

\#\# Question

To sell her plan to the city council, the architect wants to build a model to standard scale, which is one-twelfth actual size. How fast should the water flow in the channel in the model?

\#\# Answer choices

- A: A: \textbackslash{}( 0.325 \textbackslash{}text\{ m/s\} \textbackslash{})

- B: B: \textbackslash{}( 0.120 \textbackslash{}text\{ m/s\} \textbackslash{})

- C: C: \textbackslash{}( 0.491 \textbackslash{}text\{ m/s\} \textbackslash{})

- D: D: \textbackslash{}( 0.212 \textbackslash{}text\{ m/s\} \textbackslash{})

\#\# Figure caption (auxiliary; use the image as primary evidence)

The image depicts a scene where water is flowing over the edge of a concrete structure, creating a waterfall into a lower body of water. A person standing on the platform observes the waterfall. The height of the waterfall is indicated by a double-headed arrow labeled "h." The water is illustrated as a smooth, continuous flow that splashes upon hitting the water below. The scene is likely part of a diagram illustrating a concept related to fluid dynamics or physics.

\#\# Dataset metadata

Kinematics; Physical Model Grounding Reasoning, Multi-Formula Reasoning

\paragraph{Recorded answer/context.}
C: C: \textbackslash{}( 0.491 \textbackslash{}text\{ m/s\} \textbackslash{})

\paragraph{Figure/image path.}
/root/proposal\_for\_physic/science-mango/phyx\_mini\_run/phyx\_data/test\_image/684.png

\paragraph{Formalization target.}
create a compiling Lean file with sorry bodies at `PhyXMiniProblems/problem\_phyx\_mini\_0684.lean`. The Lean declarations must preserve the physical quantities, dimensions or dimensional roles, figure labels, governing-law hypotheses, and final relation expressed by this problem.
Use Mathlib/Physlib names found through LeanExplore where available. If a physics API is missing, introduce faithful local abstractions rather than scalar placeholder aliases.

\begin{theorem}[Physics formalization target]
\label{thm:physics:phyx_mini_0684:target}
\lean{PhyXMiniProblems.ProblemPhyXMini0684.problem_phyx_mini_0684}
\uses{def:physics:phyx-mini-0684:phyxminiproblems-problemphyxmini0684-speedinmeterspersecond, def:physics:phyx-mini-0684:phyxminiproblems-problemphyxmini0684-waterfallscalesetup, def:physics:phyx-mini-0684:phyxminiproblems-problemphyxmini0684-matchesprimaryfigure, def:physics:phyx-mini-0684:phyxminiproblems-problemphyxmini0684-matchesproblemdescription, def:physics:phyx-mini-0684:phyxminiproblems-problemphyxmini0684-hasphysicalparameters, def:physics:phyx-mini-0684:phyxminiproblems-problemphyxmini0684-satisfieshorizontalwaterfallkinematics, def:physics:phyx-mini-0684:phyxminiproblems-problemphyxmini0684-satisfiesstandardscalegeometry, def:physics:phyx-mini-0684:phyxminiproblems-problemphyxmini0684-answerchoice, def:physics:phyx-mini-0684:phyxminiproblems-problemphyxmini0684-answerchoice-speedinmeterspersecond, def:physics:phyx-mini-0684:phyxminiproblems-problemphyxmini0684-roundstothreedecimalplaces}
The assigned autoformalize agent should translate this physics problem into Lean declarations in the covered file, with theorem and lemma proof bodies written as `by sorry`.
\end{theorem}
\begin{proof}
This is an autoformalization task, not a proof task. Produce statements that can later be proved by the physics prover without weakening the physical model.
\end{proof}

% --- Archon declaration topology begin ---
\section{Declaration topology}
\begin{definition}[acceleration Dimension]
  \label{def:physics:phyx-mini-0684:phyxminiproblems-problemphyxmini0684-accelerationdimension}
  \lean{PhyXMiniProblems.ProblemPhyXMini0684.accelerationDimension}
  Acceleration has physical dimension L T⁻²
\end{definition}
\begin{proof}
  This is the declared model definition of acceleration Dimension.
\end{proof}
\begin{definition}[Length Quantity]
  \label{def:physics:phyx-mini-0684:phyxminiproblems-problemphyxmini0684-lengthquantity}
  \lean{PhyXMiniProblems.ProblemPhyXMini0684.LengthQuantity}
  A nonnegative, unit-independent physical length
\end{definition}
\begin{proof}
  This is the declared model definition of Length Quantity.
\end{proof}
\begin{definition}[Time Quantity]
  \label{def:physics:phyx-mini-0684:phyxminiproblems-problemphyxmini0684-timequantity}
  \lean{PhyXMiniProblems.ProblemPhyXMini0684.TimeQuantity}
  A nonnegative, unit-independent physical duration
\end{definition}
\begin{proof}
  This is the declared model definition of Time Quantity.
\end{proof}
\begin{definition}[Speed Quantity]
  \label{def:physics:phyx-mini-0684:phyxminiproblems-problemphyxmini0684-speedquantity}
  \lean{PhyXMiniProblems.ProblemPhyXMini0684.SpeedQuantity}
  A nonnegative, unit-independent speed magnitude
\end{definition}
\begin{proof}
  This is the declared model definition of Speed Quantity.
\end{proof}
\begin{definition}[Acceleration Quantity]
  \label{def:physics:phyx-mini-0684:phyxminiproblems-problemphyxmini0684-accelerationquantity}
  \lean{PhyXMiniProblems.ProblemPhyXMini0684.AccelerationQuantity}
  \uses{def:physics:phyx-mini-0684:phyxminiproblems-problemphyxmini0684-accelerationdimension}
  A nonnegative, unit-independent acceleration magnitude
\end{definition}
\begin{proof}
  This is the declared model definition of Acceleration Quantity.
\end{proof}
\begin{definition}[nonnegative SIReadout]
  \label{def:physics:phyx-mini-0684:phyxminiproblems-problemphyxmini0684-nonnegativesireadout}
  \lean{PhyXMiniProblems.ProblemPhyXMini0684.nonnegativeSIReadout}
  Read a nonnegative dimensionful quantity in coherent SI units
\end{definition}
\begin{proof}
  This is the declared model definition of nonnegative SIReadout.
\end{proof}
\begin{definition}[length In Meters]
  \label{def:physics:phyx-mini-0684:phyxminiproblems-problemphyxmini0684-lengthinmeters}
  \lean{PhyXMiniProblems.ProblemPhyXMini0684.lengthInMeters}
  \uses{def:physics:phyx-mini-0684:phyxminiproblems-problemphyxmini0684-lengthquantity, def:physics:phyx-mini-0684:phyxminiproblems-problemphyxmini0684-nonnegativesireadout}
  Metre readout of a physical length
\end{definition}
\begin{proof}
  This is the declared model definition of length In Meters.
\end{proof}
\begin{definition}[time In Seconds]
  \label{def:physics:phyx-mini-0684:phyxminiproblems-problemphyxmini0684-timeinseconds}
  \lean{PhyXMiniProblems.ProblemPhyXMini0684.timeInSeconds}
  \uses{def:physics:phyx-mini-0684:phyxminiproblems-problemphyxmini0684-timequantity, def:physics:phyx-mini-0684:phyxminiproblems-problemphyxmini0684-nonnegativesireadout}
  Second readout of a physical duration
\end{definition}
\begin{proof}
  This is the declared model definition of time In Seconds.
\end{proof}
\begin{definition}[speed In Meters Per Second]
  \label{def:physics:phyx-mini-0684:phyxminiproblems-problemphyxmini0684-speedinmeterspersecond}
  \lean{PhyXMiniProblems.ProblemPhyXMini0684.speedInMetersPerSecond}
  \uses{def:physics:phyx-mini-0684:phyxminiproblems-problemphyxmini0684-speedquantity, def:physics:phyx-mini-0684:phyxminiproblems-problemphyxmini0684-nonnegativesireadout}
  Metre-per-second readout of a speed magnitude
\end{definition}
\begin{proof}
  This is the declared model definition of speed In Meters Per Second.
\end{proof}
\begin{definition}[acceleration In Meters Per Second Squared]
  \label{def:physics:phyx-mini-0684:phyxminiproblems-problemphyxmini0684-accelerationinmeterspersecondsquared}
  \lean{PhyXMiniProblems.ProblemPhyXMini0684.accelerationInMetersPerSecondSquared}
  \uses{def:physics:phyx-mini-0684:phyxminiproblems-problemphyxmini0684-accelerationquantity, def:physics:phyx-mini-0684:phyxminiproblems-problemphyxmini0684-nonnegativesireadout}
  Metre-per-second-squared readout of an acceleration magnitude
\end{definition}
\begin{proof}
  This is the declared model definition of acceleration In Meters Per Second Squared.
\end{proof}
\begin{definition}[Channel Orientation]
  \label{def:physics:phyx-mini-0684:phyxminiproblems-problemphyxmini0684-channelorientation}
  \lean{PhyXMiniProblems.ProblemPhyXMini0684.ChannelOrientation}
  Orientation of a channel at the point where its water leaves the lip
\end{definition}
\begin{proof}
  This is the declared model definition of Channel Orientation.
\end{proof}
\begin{definition}[Receiving Region]
  \label{def:physics:phyx-mini-0684:phyxminiproblems-problemphyxmini0684-receivingregion}
  \lean{PhyXMiniProblems.ProblemPhyXMini0684.ReceivingRegion}
  Region into which the falling water is intended to land
\end{definition}
\begin{proof}
  This is the declared model definition of Receiving Region.
\end{proof}
\begin{definition}[Figure Height Label]
  \label{def:physics:phyx-mini-0684:phyxminiproblems-problemphyxmini0684-figureheightlabel}
  \lean{PhyXMiniProblems.ProblemPhyXMini0684.FigureHeightLabel}
  Literal height label visible beside the double-headed arrow in the figure
\end{definition}
\begin{proof}
  This is the declared model definition of Figure Height Label.
\end{proof}
\begin{definition}[Waterfall Figure]
  \label{def:physics:phyx-mini-0684:phyxminiproblems-problemphyxmini0684-waterfallfigure}
  \lean{PhyXMiniProblems.ProblemPhyXMini0684.WaterfallFigure}
  \uses{def:physics:phyx-mini-0684:phyxminiproblems-problemphyxmini0684-lengthquantity, def:physics:phyx-mini-0684:phyxminiproblems-problemphyxmini0684-figureheightlabel}
  The physical information transcribed from the supplied raster. The figure shows a horizontal upper channel, its lip at the top of a vertical wall, the lower pool, the curved water path, and a vertical double-headed arrow labelled h. The labelled height is itself a dimensionful length
\end{definition}
\begin{proof}
  This is the declared model definition of Waterfall Figure.
\end{proof}
\begin{definition}[Waterfall Realization]
  \label{def:physics:phyx-mini-0684:phyxminiproblems-problemphyxmini0684-waterfallrealization}
  \lean{PhyXMiniProblems.ProblemPhyXMini0684.WaterfallRealization}
  \uses{def:physics:phyx-mini-0684:phyxminiproblems-problemphyxmini0684-lengthquantity, def:physics:phyx-mini-0684:phyxminiproblems-problemphyxmini0684-timequantity, def:physics:phyx-mini-0684:phyxminiproblems-problemphyxmini0684-speedquantity, def:physics:phyx-mini-0684:phyxminiproblems-problemphyxmini0684-channelorientation, def:physics:phyx-mini-0684:phyxminiproblems-problemphyxmini0684-receivingregion}
  A physical realization of the waterfall. Fall time and horizontal landing range are independent quantities used by the kinematic laws; neither is fixed by definition from the requested speed
\end{definition}
\begin{proof}
  This is the declared model definition of Waterfall Realization.
\end{proof}
\begin{definition}[Waterfall Scale Setup]
  \label{def:physics:phyx-mini-0684:phyxminiproblems-problemphyxmini0684-waterfallscalesetup}
  \lean{PhyXMiniProblems.ProblemPhyXMini0684.WaterfallScaleSetup}
  \uses{def:physics:phyx-mini-0684:phyxminiproblems-problemphyxmini0684-accelerationquantity, def:physics:phyx-mini-0684:phyxminiproblems-problemphyxmini0684-waterfallfigure, def:physics:phyx-mini-0684:phyxminiproblems-problemphyxmini0684-waterfallrealization}
  The actual installation and its scale model experience the same local gravitational acceleration. lengthScaleFactor is dimensionless
\end{definition}
\begin{proof}
  This is the declared model definition of Waterfall Scale Setup.
\end{proof}
\begin{definition}[Matches Primary Figure]
  \label{def:physics:phyx-mini-0684:phyxminiproblems-problemphyxmini0684-matchesprimaryfigure}
  \lean{PhyXMiniProblems.ProblemPhyXMini0684.MatchesPrimaryFigure}
  \uses{def:physics:phyx-mini-0684:phyxminiproblems-problemphyxmini0684-waterfallscalesetup}
  Primary-image evidence. These fields transcribe objects and the label shown in the raster and identify its labelled height with the actual wall height; they impose no requested model-speed value
\end{definition}
\begin{proof}
  This is the declared model definition of Matches Primary Figure.
\end{proof}
\begin{definition}[Matches Problem Description]
  \label{def:physics:phyx-mini-0684:phyxminiproblems-problemphyxmini0684-matchesproblemdescription}
  \lean{PhyXMiniProblems.ProblemPhyXMini0684.MatchesProblemDescription}
  \uses{def:physics:phyx-mini-0684:phyxminiproblems-problemphyxmini0684-lengthinmeters, def:physics:phyx-mini-0684:phyxminiproblems-problemphyxmini0684-speedinmeterspersecond, def:physics:phyx-mini-0684:phyxminiproblems-problemphyxmini0684-waterfallscalesetup}
  Numerical and qualitative data stated in the problem. The actual channel speed is 1.70 m/s, the arrow height is 2.35 m, and the model is described as one-twelfth actual size. No model-speed readout occurs here
\end{definition}
\begin{proof}
  This is the declared model definition of Matches Problem Description.
\end{proof}
\begin{definition}[Has Physical Parameters]
  \label{def:physics:phyx-mini-0684:phyxminiproblems-problemphyxmini0684-hasphysicalparameters}
  \lean{PhyXMiniProblems.ProblemPhyXMini0684.HasPhysicalParameters}
  \uses{def:physics:phyx-mini-0684:phyxminiproblems-problemphyxmini0684-lengthinmeters, def:physics:phyx-mini-0684:phyxminiproblems-problemphyxmini0684-timeinseconds, def:physics:phyx-mini-0684:phyxminiproblems-problemphyxmini0684-accelerationinmeterspersecondsquared, def:physics:phyx-mini-0684:phyxminiproblems-problemphyxmini0684-waterfallscalesetup}
  Positivity and nondegeneracy conditions for the two physical falls. Although the underlying Physlib quantities are nonnegative, strict positivity is recorded where later cancellation and square-root reasoning require it
\end{definition}
\begin{proof}
  This is the declared model definition of Has Physical Parameters.
\end{proof}
\begin{definition}[Satisfies Horizontal Waterfall Kinematics]
  \label{def:physics:phyx-mini-0684:phyxminiproblems-problemphyxmini0684-satisfieshorizontalwaterfallkinematics}
  \lean{PhyXMiniProblems.ProblemPhyXMini0684.SatisfiesHorizontalWaterfallKinematics}
  \uses{def:physics:phyx-mini-0684:phyxminiproblems-problemphyxmini0684-lengthinmeters, def:physics:phyx-mini-0684:phyxminiproblems-problemphyxmini0684-timeinseconds, def:physics:phyx-mini-0684:phyxminiproblems-problemphyxmini0684-speedinmeterspersecond, def:physics:phyx-mini-0684:phyxminiproblems-problemphyxmini0684-accelerationinmeterspersecondsquared, def:physics:phyx-mini-0684:phyxminiproblems-problemphyxmini0684-waterfallscalesetup}
  For each realization, water has zero initial vertical speed, undergoes constant-gravity vertical free fall, and retains its horizontal channel speed. The relations are written in coherent SI readouts so their scalar equations are dimensionally homogeneous. They contain no numerical model speed
\end{definition}
\begin{proof}
  This is the declared model definition of Satisfies Horizontal Waterfall Kinematics.
\end{proof}
\begin{definition}[Satisfies Standard Scale Geometry]
  \label{def:physics:phyx-mini-0684:phyxminiproblems-problemphyxmini0684-satisfiesstandardscalegeometry}
  \lean{PhyXMiniProblems.ProblemPhyXMini0684.SatisfiesStandardScaleGeometry}
  \uses{def:physics:phyx-mini-0684:phyxminiproblems-problemphyxmini0684-lengthinmeters, def:physics:phyx-mini-0684:phyxminiproblems-problemphyxmini0684-waterfallscalesetup}
  Standard geometric similarity scales every relevant length by the common dimensionless factor and preserves the qualitative channel/pool geometry. Scaling the landing range, rather than merely the wall height, records that the entire water trajectory is to have the same shape in the model
\end{definition}
\begin{proof}
  This is the declared model definition of Satisfies Standard Scale Geometry.
\end{proof}
\begin{definition}[Answer Choice]
  \label{def:physics:phyx-mini-0684:phyxminiproblems-problemphyxmini0684-answerchoice}
  \lean{PhyXMiniProblems.ProblemPhyXMini0684.AnswerChoice}
  Labels of the four answer choices printed with the problem
\end{definition}
\begin{proof}
  This is the declared model definition of Answer Choice.
\end{proof}
\begin{definition}[speed In Meters Per Second]
  \label{def:physics:phyx-mini-0684:phyxminiproblems-problemphyxmini0684-answerchoice-speedinmeterspersecond}
  \lean{PhyXMiniProblems.ProblemPhyXMini0684.AnswerChoice.speedInMetersPerSecond}
  \uses{def:physics:phyx-mini-0684:phyxminiproblems-problemphyxmini0684-speedinmeterspersecond, def:physics:phyx-mini-0684:phyxminiproblems-problemphyxmini0684-answerchoice}
  Metre-per-second value printed beside each displayed answer choice
\end{definition}
\begin{proof}
  This is the declared model definition of speed In Meters Per Second.
\end{proof}
\begin{definition}[Rounds To Three Decimal Places]
  \label{def:physics:phyx-mini-0684:phyxminiproblems-problemphyxmini0684-roundstothreedecimalplaces}
  \lean{PhyXMiniProblems.ProblemPhyXMini0684.RoundsToThreeDecimalPlaces}
  Rounding to the nearest 0.001 is represented by an error strictly less than half of that unit. The exact ideal-model speed is not exactly the decimal 0.491, so making this precision convention explicit avoids a false equality
\end{definition}
\begin{proof}
  This is the declared model definition of Rounds To Three Decimal Places.
\end{proof}
% --- Archon declaration topology end ---
% --- Archon physics formalization source end ---
